\chapter{Theophylline Level}
\section*{What Is This Test?}
The theophylline level measures the serum concentration of theophylline, a methylxanthine bronchodilator used to treat asthma and chronic obstructive pulmonary disease (COPD). Theophylline causes bronchodilation by inhibiting phosphodiesterase (PDE3 and PDE4), increasing intracellular cyclic AMP (cAMP) in bronchial smooth muscle, and by antagonizing adenosine receptors (A1 and A2). Theophylline use has declined significantly with the advent of inhaled beta-agonists and inhaled corticosteroids; however, it is still used as an add-on therapy for difficult-to-control asthma and COPD. Theophylline has a narrow therapeutic index (10--20 mcg/mL). Toxicity occurs at serum concentrations >20 mcg/mL: gastrointestinal (nausea, vomiting, abdominal pain, hematemesis from gastritis), cardiac (sinus tachycardia, supraventricular arrhythmias, ventricular arrhythmias), and neurologic (tremor, agitation, seizures --- seizures are a life-threatening manifestation of severe theophylline toxicity and can occur without preceding mild symptoms, particularly in acute overdose). Theophylline is metabolized by CYP1A2 in the liver. Drugs that inhibit CYP1A2 (ciprofloxacin, fluvoxamine, cimetidine) and conditions that impair metabolism (congestive heart failure, liver disease, advanced age) increase theophylline levels and predispose to toxicity.
\section*{Alternative Names}
Theophylline Level; Aminophylline Level; Theodur Level; Serum Theophylline; Methylxanthine Level
\section*{Principle}
Theophylline is measured by particle-enhanced turbidimetric inhibition immunoassay (PETINIA) or enzyme-multiplied immunoassay technique (EMIT) on automated analyzers. Caffeine (structurally similar methylxanthine) cross-reacts with some immunoassays and produces falsely elevated theophylline levels.
\section*{History}
Theophylline is a naturally occurring methylxanthine present in tea (Camellia sinensis) that was first isolated from tea leaves in 1888 by Albrecht Kossel. Its bronchodilator properties were described in the early 20th century, and aminophylline (a theophylline-ethylenediamine salt) was introduced as a parenteral bronchodilator for acute asthma in 1937. Sustained-release oral preparations introduced in the 1970s--1980s made chronic oral therapy practical, and serum immunoassays enabled routine TDM of this narrow-therapeutic-index drug. Although theophylline was once a first-line asthma therapy, inhaled corticosteroids and long-acting beta-agonists have largely replaced it, and it is now used mainly as an add-on for difficult-to-control asthma and COPD.
\section*{Why This Test Is Ordered}
TDM to maintain levels in the therapeutic range. Diagnosis of theophylline toxicity. Dose adjustment in patients with altered CYP1A2 metabolism.
\section*{Primary vs Secondary Test}
TDM is a secondary monitoring test. Theophylline therapy is primarily guided by clinical response (symptom control, spirometry, peak flow) and side-effect tolerance; the serum level is a secondary tool used to confirm that dosing is therapeutic, to detect or avoid toxicity given the drug's narrow therapeutic index, and to guide dose adjustments when metabolism-altering drugs or comorbidities are present.
\section*{How This Test Is Performed}
A blood sample is collected by standard venipuncture into a serum separator tube (SST). Theophylline is measured in serum by immunoassay (PETINIA or EMIT) on automated analyzers, with a typical turnaround time of 1--2 hours for urgent requests and 1--2 days for routine samples. LC-MS/MS is available in some centers as a reference method without caffeine cross-reactivity.
\section*{What Preparation Is Needed}
The sample should be drawn as a trough, immediately before the next scheduled dose of a sustained-release oral preparation, at steady state (approximately 2--3 days after initiation or dose change). Patients should take theophylline at consistent times each day, and the dose, formulation, and time of the last dose should be documented. Because caffeine (coffee, tea, cola, chocolate) cross-reacts with some immunoassays, large caffeinated intake immediately before sampling should be avoided. A current medication list is essential, since CYP1A2 inhibitors and inducers markedly alter the level.
\section*{Contraindications and Precautions}
\begin{itemize}
  \item No absolute contraindications to standard venipuncture
  \item Suspected or established theophylline toxicity (severe nausea, vomiting, tachycardia, tremor, agitation, seizures) is a clinical emergency that must be treated on the basis of symptoms and history without waiting for the level
  \item Interpretation requires knowledge of concurrent drugs: CYP1A2 inhibitors (ciprofloxacin, fluvoxamine, cimetidine, oral contraceptives) increase levels; CYP1A2 inducers (smoking, rifampin, phenytoin, carbamazepine) decrease levels
  \item Conditions that slow clearance (congestive heart failure, cirrhosis, fever, advanced age) lower the toxic threshold, and in acute overdose severe toxicity can occur at levels well tolerated in chronic therapy
  \item Caffeine ingestion and immunoassay cross-reactivity can falsely elevate the measured level; interpret in clinical context
\end{itemize}
\section*{Risks}
Minimal risk. Slight pain or bruising at the venipuncture site. Rarely: hematoma, infection, or vasovagal syncope.
\section*{What Happens After the Test}
The result is interpreted against the therapeutic range (10--20 mcg/mL; 5--15 mcg/mL in some guidelines) in the context of symptoms, comorbidities, and interacting drugs. Subtherapeutic levels prompt a dose increase with repeat sampling after a new steady state (2--3 days); toxic levels prompt dose reduction or discontinuation, supportive care, and serial levels until the concentration falls into range. A follow-up level is typically obtained 2--3 days after any dose change.
\section*{Understanding Results}
Subtherapeutic (<10 mcg/mL): Inadequate bronchodilation, breakthrough asthma symptoms. Toxic (>20 mcg/mL): nausea/vomiting (common early sign), tachycardia, tremor, agitation. Severe toxicity (>40 mcg/mL): cardiac arrhythmias (supraventricular tachycardia, atrial fibrillation, ventricular tachycardia), seizures, cardiovascular collapse. Seizures from theophylline toxicity are difficult to control and associated with high mortality.
\section*{Normal Results}
Therapeutic range: 10--20 mcg/mL for bronchodilation. Some guidelines recommend 5--15 mcg/mL as adequate with fewer side effects. Trough: drawn just before the next scheduled dose (for sustained-release oral formulations). For IV aminophylline continuous infusion: steady state level at any time during the infusion.
\section*{Follow-up Tests}
ECG (tachycardia, arrhythmias), serum electrolytes (hypokalemia, hypomagnesemia from beta-adrenergic stimulation), serum glucose (hyperglycemia from catecholamine release), liver function tests, and activated charcoal for acute overdose.
\section*{References}
\begin{enumerate}
  \item Barnes PJ. Theophylline. Am J Respir Crit Care Med. 2013;188(8):901--906.
  \item Hendeles L, Weinberger M. Theophylline: a state of the art review. Pharmacotherapy. 1983;3(1):2--44.
  \item Weinberger M, Hendeles L. Theophylline in asthma. N Engl J Med. 1996;334(21):1380--1388.
  \item Mitenko PA, Ogilvie RI. Rational intravenous doses of theophylline. N Engl J Med. 1973;289(12):600--603.
  \item Burtis CA, Ashwood ER, Bruns DE. Tietz Textbook of Clinical Chemistry and Molecular Diagnostics. 5th ed. St. Louis, MO: Elsevier Saunders; 2012.
\end{enumerate}
\clearpage
\section*{Regulatory Status and Authorization Date}
Regulatory status applies to the specific assay, instrument, manufacturer, and intended use rather than to the test name alone. No single approval date can be assigned to this test category. For a commercial assay, verify the current product in the relevant regulator's database: the U.S. Food and Drug Administration (FDA) clearance, approval, or authorization; India's Central Drugs Standard Control Organisation (CDSCO) licence or permission; the European Union In Vitro Diagnostic Regulation (EU IVDR) CE certificate issued through the applicable conformity-assessment route; or the United Kingdom Medicines and Healthcare products Regulatory Agency (MHRA) registration and UKCA/CE status. Laboratory-developed tests may not have an individual FDA, CDSCO, EU IVDR, or MHRA approval date.
\clearpage
